\section{2019 Algebra & Number Theory}

\subsection*{Individual}
\begin{exercise}
Let $G$ be a finite group. Assume that for any representation $V$ of $G$ over a field of characteristic zero, the character $\chi_V$ takes value in $\mathbb{Q}$. Assume $g$ is an element in $G$ such that $g^{2019} = 1$.

Prove that $g$ and $g^{19}$ are conjugate in $G$.
\end{exercise}

\begin{solution}
A well-known result in the representation theory of finite groups states that all characters of a finite group $G$ are rational-valued if and only if for every element $g \in G$, $g$ is conjugate to $g^k$ for all integers $k$ such that $\gcd(k, |g|) = 1$.

In this problem, we are given that for any representation $V$ over a field of characteristic zero, the character $\chi_V$ takes values in $\mathbb{Q}$. This implies that for any $g \in G$, $g$ is conjugate to $g^k$ for all $k$ coprime to the order of $g$.

We are given $g^{2019} = 1$, so the order of $g$, denoted by $|g|$, must be a divisor of $2019$. The prime factorization of $2019$ is $3 \times 673$. Thus, $|g| \in \{1, 3, 673, 2019\}$.
To show that $g$ and $g^{19}$ are conjugate, it suffices to verify that $\gcd(19, |g|) = 1$. Since $19$ is a prime number and $19$ does not divide $3$ or $673$, it follows that $19$ is coprime to any divisor of $2019$ (except possibly if the divisor is a multiple of $19$, which is not the case here).

Therefore, $\gcd(19, |g|) = 1$ for all possible orders $|g|$, and by the property of rational characters, $g$ and $g^{19}$ are conjugate in $G$.
\end{solution}

\begin{exercise}
Let $p$ be a prime number, and let $\mathbb{F}_p$ be the finite field with $p$ elements. Let $F = \mathbb{F}_p(t)$ be the field of rational functions over $\mathbb{F}_p$. Consider all subfields $C$ of $F$ such that $F/C$ is a finite Galois extension.

1. Show that among such subfields, there is a smallest one $C_0$, i.e., $C_0$ is contained in any other $C$.

2. What is the degree of $F/C_0$?
\end{exercise}

\begin{solution}
1. By the Fundamental Theorem of Galois Theory for function fields, if $F/C$ is a finite Galois extension, then $C = F^G$ for some finite subgroup $G$ of the automorphism group $\mathrm{Aut}(F)$.
The automorphism group of the field of rational functions $F = \mathbb{F}_p(t)$ is the group of fractional linear transformations:
$$\mathrm{Aut}(F) \cong \mathrm{PGL}_2(\mathbb{F}_p).$$
Since $\mathbb{F}_p$ is a finite field, $\mathrm{PGL}_2(\mathbb{F}_p)$ is a finite group. Every subfield $C$ such that $F/C$ is a finite Galois extension corresponds to a subgroup $G \subset \mathrm{PGL}_2(\mathbb{F}_p)$.
The inclusion of subfields $C_1 \subset C_2$ corresponds to the inclusion of subgroups $G_2 \subset G_1$. Therefore, the smallest subfield $C_0$ corresponds to the largest possible finite subgroup of $\mathrm{Aut}(F)$.
Since $\mathrm{Aut}(F) = \mathrm{PGL}_2(\mathbb{F}_p)$ is itself finite, the largest subgroup is $G_0 = \mathrm{PGL}_2(\mathbb{F}_p)$.
Thus, $C_0 = F^{\mathrm{PGL}_2(\mathbb{F}_p)}$ is the smallest subfield. Any other such $C$ corresponds to a subgroup $G \subseteq \mathrm{PGL}_2(\mathbb{F}_p)$, so $C = F^G \supseteq F^{\mathrm{PGL}_2(\mathbb{F}_p)} = C_0$.

2. The degree of the extension $F/C_0$ is equal to the order of the Galois group $G_0 = \mathrm{PGL}_2(\mathbb{F}_p)$.
The order of $\mathrm{GL}_2(\mathbb{F}_p)$ is $(p^2-1)(p^2-p)$.
The group $\mathrm{PGL}_2(\mathbb{F}_p)$ is the quotient of $\mathrm{GL}_2(\mathbb{F}_p)$ by its center (scalar matrices), which has order $p-1$.
Thus,
$$[F : C_0] = |\mathrm{PGL}_2(\mathbb{F}_p)| = \frac{(p^2-1)(p^2-p)}{p-1} = (p+1)p(p-1) = p^3 - p.$$
\end{solution}

\begin{exercise}
Let $R \subset R'$ be an integral extension of commutative rings. Let $\mathfrak{p}'$ be a prime ideal of $R'$. Prove that $\mathfrak{p}'$ is a maximal ideal of $R'$ if and only if $\mathfrak{p}' \cap R$ is a maximal ideal of $R$.
\end{exercise}

\begin{solution}
Let $\mathfrak{p} = \mathfrak{p}' \cap R$. Since $\mathfrak{p}'$ is a prime ideal of $R'$, its intersection with the subring $R$ is a prime ideal of $R$. The inclusion $R \subseteq R'$ induces an injective homomorphism of integral domains:
$$ \phi: R/\mathfrak{p} \hookrightarrow R'/\mathfrak{p}' $$
Since $R'$ is integral over $R$, the quotient $R'/\mathfrak{p}'$ is integral over $R/\mathfrak{p}$. We use the following standard result in commutative algebra:
\textit{Let $D \subseteq D'$ be an integral extension of integral domains. Then $D'$ is a field if and only if $D$ is a field.}

\begin{itemize}
    \item $(\implies)$ Assume $\mathfrak{p}'$ is maximal. Then $R'/\mathfrak{p}'$ is a field. Since $R'/\mathfrak{p}'$ is integral over $R/\mathfrak{p}$, and $R'/\mathfrak{p}'$ is a field, then $R/\mathfrak{p}$ must be a field. Indeed, for any non-zero $x \in R/\mathfrak{p}$, $x^{-1} \in R'/\mathfrak{p}'$ exists. Since $x^{-1}$ is integral over $R/\mathfrak{p}$, there is an equation
    $$ (x^{-1})^n + a_{n-1}(x^{-1})^{n-1} + \cdots + a_0 = 0, \quad a_i \in R/\mathfrak{p}. $$
    Multiplying by $x^{n-1}$, we get $x^{-1} = -(a_{n-1} + a_{n-2}x + \cdots + a_0 x^{n-1}) \in R/\mathfrak{p}$. Thus $R/\mathfrak{p}$ is a field, so $\mathfrak{p}$ is maximal in $R$.

    \item $(\impliedby)$ Assume $\mathfrak{p} = \mathfrak{p}' \cap R$ is maximal. Then $R/\mathfrak{p}$ is a field. Since $R'/\mathfrak{p}'$ is an integral domain that is integral over a field $R/\mathfrak{p}$, it must be a field. For any non-zero $\alpha \in R'/\mathfrak{p}'$, $\alpha$ satisfies an equation
    $$ \alpha^m + c_{m-1}\alpha^{m-1} + \cdots + c_1 \alpha + c_0 = 0, \quad c_i \in R/\mathfrak{p}. $$
    Since $R'/\mathfrak{p}'$ is a domain, we can assume $c_0 \neq 0$ (otherwise divide by $\alpha$). Then
    $$ \alpha(\alpha^{m-1} + \cdots + c_1) = -c_0 \implies \alpha \cdot \left[ -c_0^{-1}(\alpha^{m-1} + \cdots + c_1) \right] = 1. $$
    Thus $\alpha$ is invertible, so $R'/\mathfrak{p}'$ is a field, meaning $\mathfrak{p}'$ is maximal in $R'$.
\end{itemize}
\end{solution}

\begin{exercise}
1. Prove that $\mathrm{GL}_n(\mathbb{C})$ is path-connected.

2. Let $X = \{A \in \mathrm{GL}_n(\mathbb{C}) \mid A^m = \mathrm{Id}\}$. Describe the path-connected component of $X$ and prove your answer.
\end{exercise}

\begin{solution}
1. $\mathrm{GL}_n(\mathbb{C})$ is path-connected. This can be seen in several ways:
\begin{itemize}
    \item \textbf{Topological Argument:} $\mathrm{GL}_n(\mathbb{C})$ is an open subset of $M_n(\mathbb{C}) \cong \mathbb{C}^{n^2}$. The complement is the zero set of the polynomial $\det(A)$. In a complex space, the zero set of a non-constant polynomial has complex codimension 1 (real codimension 2), and thus its complement is path-connected.
    \item \textbf{Explicit Path:} For any $A \in \mathrm{GL}_n(\mathbb{C})$, there exists a Jordan normal form $J = P A P^{-1}$. We can connect $J$ to a diagonal matrix by scaling off-diagonal entries to zero. For a diagonal matrix $\mathrm{diag}(\lambda_1, \ldots, \lambda_n)$, we can connect each $\lambda_i \in \mathbb{C}^*$ to $1$ via a path in $\mathbb{C}^*$ (which is path-connected).
\end{itemize}

2. Let $\omega = e^{2\pi i / m}$ be a primitive $m$-th root of unity. Any $A \in X$ satisfies $A^m - I = 0$. Since the polynomial $z^m - 1$ has distinct roots $\{1, \omega, \ldots, \omega^{m-1}\}$, $A$ is diagonalizable and its eigenvalues must be in this set.
Let $n_k$ be the multiplicity of $\omega^k$ as an eigenvalue of $A$. Then $\sum_{k=0}^{m-1} n_k = n$.
The set $X$ is the union of conjugacy orbits in $\mathrm{GL}_n(\mathbb{C})$. Two matrices in $X$ are conjugate if and only if they have the same eigenvalues with the same multiplicities. Let $\mathcal{O}_{(n_0, \ldots, n_{m-1})}$ be the conjugacy orbit of the diagonal matrix with $n_k$ entries of $\omega^k$.
Each orbit $\mathcal{O}_{\mathbf{n}}$ is the image of the continuous map $\mathrm{GL}_n(\mathbb{C}) \to \mathrm{GL}_n(\mathbb{C})$ given by $P \mapsto P D_{\mathbf{n}} P^{-1}$. Since $\mathrm{GL}_n(\mathbb{C})$ is connected, each orbit is a connected set.
Moreover, the multiplicities of the eigenvalues are discrete data. The map from $X$ to the set of multisets of eigenvalues (or the coefficients of the characteristic polynomial) is continuous. Since the roots are fixed (they are $m$-th roots of unity), the only way for the roots to change is to jump from one $\omega^j$ to another, which is impossible in a connected component.
Thus, the path-connected components of $X$ are precisely the conjugacy orbits $\mathcal{O}_{(n_0, \ldots, n_{m-1})}$. The number of such components is the number of ways to write $n$ as a sum of $m$ non-negative integers, which is $\binom{n+m-1}{m-1}$.
\end{solution}

\begin{exercise}
The Fibonacci sequence is defined by
$$F_0 = 0, \quad F_1 = 1, \quad F_{n+2} = F_{n+1} + F_n.$$
Let $p$ be a prime number.

1. Show that if $p \equiv 1, 4 \pmod{5}$, then $p$ divides $F_{p-1}$.

2. Let $\mathbb{F}_{p^2}$ be the finite field of $p^2$ elements. Show that the norm map $N: \mathbb{F}_{p^2}^\times \to \mathbb{F}_p^\times$ is surjective, and deduce the cardinality of the kernel of $N$.

3. Show that if $p \equiv 2, 3 \pmod{5}$, then $p$ divides $F_{p+1}$.
\end{exercise}

\begin{solution}
The Fibonacci sequence is given by Binet's formula: $F_n = \frac{\alpha^n - \beta^n}{\alpha - \beta}$, where $\alpha, \beta$ are the roots of $x^2 - x - 1 = 0$, namely $\frac{1 \pm \sqrt{5}}{2}$. The discriminant of this quadratic is $5$.

1. If $p \equiv 1, 4 \pmod 5$, then by Quadratic Reciprocity, $\left(\frac{5}{p}\right) = \left(\frac{p}{5}\right) = 1$. Thus $\sqrt{5}$ exists in $\mathbb{F}_p$, and $\alpha, \beta \in \mathbb{F}_p$. By Fermat's Little Theorem, for $p \neq 2, 5$, we have $\alpha^{p-1} \equiv 1 \pmod p$ and $\beta^{p-1} \equiv 1 \pmod p$.
Thus $F_{p-1} = \frac{\alpha^{p-1} - \beta^{p-1}}{\alpha - \beta} \equiv \frac{1 - 1}{\sqrt{5}} \equiv 0 \pmod p$.
(For $p=5$, $F_4 = 3$, not div by 5. For $p=2$, $F_1 = 1$. The condition $p \equiv 1, 4 \pmod 5$ excludes $2$ and $5$).

2. The norm map $N: \mathbb{F}_{p^2}^\times \to \mathbb{F}_p^\times$ is defined by $N(x) = x \cdot x^p = x^{p+1}$.
Since $\mathbb{F}_{p^2}^\times$ is a cyclic group of order $p^2 - 1$, and $x \mapsto x^{p+1}$ is a power map, the image has size $\frac{p^2-1}{\gcd(p+1, p^2-1)} = \frac{p^2-1}{p+1} = p - 1$.
Since $|\mathbb{F}_p^\times| = p - 1$, the norm map is surjective.
The kernel is a subgroup of $\mathbb{F}_{p^2}^\times$ of order $\frac{p^2 - 1}{p - 1} = p + 1$.

3. If $p \equiv 2, 3 \pmod 5$, then $\left(\frac{5}{p}\right) = \left(\frac{p}{5}\right) = -1$. Thus $\sqrt{5} \notin \mathbb{F}_p$, which implies $\alpha, \beta \in \mathbb{F}_{p^2} \setminus \mathbb{F}_p$.
The Frobenius automorphism $\sigma(x) = x^p$ interchanges the roots of $x^2 - x - 1 = 0$ because they are conjugated in the quadratic extension $\mathbb{F}_p(\sqrt{5})$. Thus $\alpha^p = \beta$ and $\beta^p = \alpha$.
The norm of $\alpha$ is $\alpha \cdot \alpha^p = \alpha \beta = -1$ (from the product of roots of $x^2 - x - 1 = 0$).
So $\alpha^{p+1} = -1$ and $\beta^{p+1} = -1$.
Then $F_{p+1} = \frac{\alpha^{p+1} - \beta^{p+1}}{\alpha - \beta} \equiv \frac{-1 - (-1)}{\sqrt{5}} \equiv 0 \pmod p$.
\end{solution}

\subsection*{Team}
\begin{exercise}
Let $S_n$ be the group of permutations of $\{1, 2, \ldots, n\}$. Let $\sigma \in S_n$ be the permutation
$$(1, n)(2, n-1)\cdots(k, n-k+1)\cdots.$$
Prove that the centralizer $Z_{S_n}(\sigma)$ is isomorphic to $S_{[n/2]} \ltimes (\mathbb{Z}/2\mathbb{Z})^{[n/2]}$.
\end{exercise}

\begin{solution}
Let $k = [n/2]$. The permutation $\sigma$ is a product of $k$ disjoint transpositions:
$$\sigma = \tau_1 \tau_2 \cdots \tau_k, \quad \text{where } \tau_j = (j, n-j+1).$$
If $n = 2k$, $\sigma$ has no fixed points. If $n = 2k+1$, $\sigma$ has exactly one fixed point $k+1$.

An element $g \in S_n$ commutes with $\sigma$ if and only if $g \sigma g^{-1} = \sigma$. Since conjugation preserves the cycle structure, $g$ must map the 2-cycles of $\sigma$ to 2-cycles of $\sigma$.
The 2-cycles of $\sigma$ are the sets $B_j = \{j, n-j+1\}$ for $j=1, \ldots, k$.
Any $g \in Z_{S_n}(\sigma)$ must permute these $k$ sets $\{B_1, \ldots, B_k\}$. This defines a homomorphism $\phi: Z_{S_n}(\sigma) \to S_k$, where $S_k$ is the symmetric group on the set of blocks $\{B_j\}$.

The kernel of $\phi$ consists of those $g$ that satisfy $g(B_j) = B_j$ for all $j$. For each $j$, $g$ can either be the identity on $B_j$ or the transposition $(j, n-j+1)$. Thus the kernel is isomorphic to $(\mathbb{Z}/2\mathbb{Z})^k$.
If $n = 2k+1$, the element $g$ must also map the fixed point $k+1$ to a fixed point of $\sigma$. Since $k+1$ is the unique fixed point, $g(k+1) = k+1$. Thus the fixed point does not add any new degrees of freedom.

The map $\phi$ is surjective. For any $\pi \in S_k$, we can define $g_\pi \in S_n$ such that $g_\pi(j) = \pi(j)$ and $g_\pi(n-j+1) = n-\pi(j)+1$. This $g_\pi$ commutes with $\sigma$ and $\phi(g_\pi) = \pi$.
This shows that $Z_{S_n}(\sigma)$ is an extension of $S_k$ by $(\mathbb{Z}/2\mathbb{Z})^k$. Since the lifting $g_\pi$ provides a section, the extension is a semidirect product:
$$Z_{S_n}(\sigma) \cong S_k \ltimes (\mathbb{Z}/2\mathbb{Z})^k.$$
The action of $S_k$ on $(\mathbb{Z}/2\mathbb{Z})^k$ is the standard action by permutation of the coordinates. This group is also known as the hyperoctahedral group or the wreath product $S_2 \wr S_k$.
\end{solution}

\begin{exercise}
Recall that the algebra of regular functions on a vector space $W$ is the symmetric algebra of linear forms on $W$.

Let $V = \mathbb{C}^2$. Let $\mathbb{C}[\mathrm{End}_{\mathbb{C}}(V)]$ be the algebra of regular functions on $\mathrm{End}_{\mathbb{C}}(V)$. The natural action of the group $G = SL_2(\mathbb{C})$ on $V$ induces an action of $G \times G$ on $\mathrm{End}_{\mathbb{C}}(V)$ by left and right multiplication. Thus we get an action of $G \times G$ on $\mathbb{C}[\mathrm{End}_{\mathbb{C}}(V)]$.

Compute the algebra of fixed points $\mathbb{C}[\mathrm{End}_{\mathbb{C}}(V)]^{G \times G}$.
\end{exercise}

\begin{solution}
Let $V = \mathbb{C}^2$ and $\mathrm{End}(V) \cong M_2(\mathbb{C})$. The group $G = SL_2(\mathbb{C})$ acts on $V$ and its dual $V^*$. The action of $G \times G$ on $X \in \mathrm{End}(V)$ is given by $(g_1, g_2) \cdot X = g_1 X g_2^{-1}$.
The algebra of regular functions is $\mathbb{C}[M_2(\mathbb{C})] \cong S^*(\mathrm{End}(V)^*)$.
As a $G \times G$-module, we have $\mathrm{End}(V) \cong V \otimes V^*$. Since $V \cong V^*$ for $SL_2(\mathbb{C})$ (via the invariant skew-symmetric form), we have $\mathrm{End}(V) \cong V \otimes V$ as a $G \times G$ module.
By the Cauchy identity for the symmetric algebra:
$$ S^d(V \otimes V) \cong \bigoplus_{\lambda \vdash d, \ell(\lambda) \leq 2} S^\lambda V \otimes S^\lambda V $$
where $S^\lambda V$ are the Schur modules (irreducible representations of $GL(V)$). For $SL_2(\mathbb{C})$, the irreducible representations are $V_k = S^k V$ of dimension $k+1$.
The Schur module $S^{(a,b)} V$ for $SL_2(\mathbb{C})$ is isomorphic to $V_{a-b}$ (the $(a-b)$-th symmetric power).
An invariant in $S^\lambda V \otimes S^\lambda V$ under $G \times G$ exists if and only if $S^\lambda V$ is the trivial representation.
For $SL_2(\mathbb{C})$, $V_{a-b}$ is trivial if and only if $a-b = 0$, i.e., $a=b$.
This implies that $\lambda = (k, k)$ for some $k \geq 0$. Then $d = 2k$.
For each $k$, there is exactly one such partition, giving a 1-dimensional space of invariants in degree $2k$.
In degree $d=2$ (where $k=1$), the partition is $(1, 1)$, and the invariant is the determinant function $\det(X)$.
Since $\det(g_1 X g_2^{-1}) = \det(g_1) \det(X) \det(g_2)^{-1} = \det(X)$ for $g_1, g_2 \in SL_2(\mathbb{C})$, the determinant is indeed an invariant.
The invariants in higher degrees $2k$ are simply the powers $\det(X)^k$.
Thus, the algebra of fixed points is:
$$ \mathbb{C}[\mathrm{End}_{\mathbb{C}}(V)]^{G \times G} = \mathbb{C}[\det]. $$
\end{solution}

\begin{exercise}
Let $R$ be a Noetherian ring and $I \subset R$ be an ideal. Define the Rees algebra as
$$\mathrm{Rees}(I, R) := \bigoplus_{n \geq 1} I^n t^n \subset R[t].$$
Prove that $\mathrm{Rees}(I, R)$ is Noetherian.
\end{exercise}

\begin{solution}
The Rees algebra is defined as $S = \mathrm{Rees}(I, R) = \bigoplus_{n=0}^\infty I^n t^n = R \oplus It \oplus I^2 t^2 \oplus \cdots \subseteq R[t]$. (Note: the problem statement had $n \geq 1$, but for it to be an algebra with identity, the $n=0$ term $I^0 = R$ is included).

Since $R$ is a Noetherian ring, every ideal is finitely generated. Let $I = (a_1, a_2, \dots, a_k)$ where $a_i \in R$.
We claim that $S$ is finitely generated as an $R$-algebra by the elements $\{a_1 t, a_2 t, \dots, a_k t\}$.
Indeed, any element in $I^n t^n$ is a sum of terms of the form $b t^n$ where $b \in I^n$. Since $b \in I^n$, $b$ is a sum of products $x_1 x_2 \cdots x_n$ where each $x_j \in \{a_1, \dots, a_k\}$.
Thus $b t^n$ is a sum of products $(x_1 t)(x_2 t) \cdots (x_n t)$.
This shows that $S = R[a_1 t, a_2 t, \dots, a_k t]$.

By the Hilbert Basis Theorem, if $R$ is a Noetherian ring, then the polynomial ring $R[x_1, \dots, x_k]$ is also Noetherian.
The Rees algebra $S$ is a homomorphic image of the polynomial ring $R[x_1, \dots, x_k]$ under the map $x_i \mapsto a_i t$.
Since any quotient of a Noetherian ring is Noetherian, it follows that $\mathrm{Rees}(I, R)$ is Noetherian.
\end{solution}

\begin{exercise}
Let $p$ be an odd prime. Let $\Phi_p$ be the $p$-th cyclotomic field, i.e., $\Phi_p = \mathbb{Q}(\zeta_p)$ where $\zeta_p$ is a primitive $p$-th root of unity.

1. Show that $\Phi_p/\mathbb{Q}$ is a Galois extension with Galois group $(\mathbb{Z}/p\mathbb{Z})^\times$.

2. Deduce that $\Phi_p$ contains a unique quadratic extension of $\mathbb{Q}$.

3. Write $g_p$ for the Gauss sum $\sum_{a=1}^{p-1} \left(\frac{a}{p}\right) \zeta_p^a$. Show that $\overline{g_p} = \left(\frac{-1}{p}\right) g_p$.

4. Determine the unique quadratic extension of $\mathbb{Q}$ contained in $\Phi_p$.
\end{exercise}

\begin{solution}
1. The minimal polynomial of $\zeta_p$ over $\mathbb{Q}$ is the $p$-th cyclotomic polynomial $\Phi_p(x) = x^{p-1} + x^{p-2} + \dots + 1$, which is irreducible of degree $p-1$. The roots are $\{\zeta_p, \zeta_p^2, \dots, \zeta_p^{p-1}\}$. Since all roots are in $\mathbb{Q}(\zeta_p)$, the extension is Galois. Any $\sigma \in \mathrm{Gal}(\Phi_p/\mathbb{Q})$ is determined by $\sigma(\zeta_p) = \zeta_p^a$ for some $a \in (\mathbb{Z}/p\mathbb{Z})^\times$. The map $\sigma \mapsto a$ is an isomorphism $\mathrm{Gal}(\Phi_p/\mathbb{Q}) \cong (\mathbb{Z}/p\mathbb{Z})^\times$.

2. The Galois group $(\mathbb{Z}/p\mathbb{Z})^\times$ is a cyclic group of order $p-1$. Since $p$ is an odd prime, $p-1$ is even. A cyclic group of even order $n$ has a unique subgroup of order $n/2$ (the subgroup of squares). By the Fundamental Theorem of Galois Theory, there is a unique subfield $K \subset \Phi_p$ such that $[K : \mathbb{Q}] = 2$, corresponding to the subgroup of squares $H = \{ a^2 \mid a \in (\mathbb{Z}/p\mathbb{Z})^\times \}$.

3. The Gauss sum is $g_p = \sum_{a=1}^{p-1} \left(\frac{a}{p}\right) \zeta_p^a$.
The complex conjugate is $\overline{g_p} = \sum_{a=1}^{p-1} \left(\frac{a}{p}\right) \zeta_p^{-a}$.
Since $\left(\frac{a}{p}\right) = \left(\frac{-1}{p}\right) \left(\frac{-a}{p}\right)$, we have:
$\overline{g_p} = \sum_{a=1}^{p-1} \left(\frac{-1}{p}\right) \left(\frac{-a}{p}\right) \zeta_p^{-a}$.
Let $b = -a \pmod p$. As $a$ runs from $1$ to $p-1$, $b$ also runs from $1$ to $p-1$.
$\overline{g_p} = \left(\frac{-1}{p}\right) \sum_{b=1}^{p-1} \left(\frac{b}{p}\right) \zeta_p^{b} = \left(\frac{-1}{p}\right) g_p$.

4. A well-known property of Gauss sums is that $g_p^2 = \left(\frac{-1}{p}\right) p$.
By part 3, if $\left(\frac{-1}{p}\right) = 1$ (i.e., $p \equiv 1 \pmod 4$), then $g_p$ is real and $g_p^2 = p$, so $\mathbb{Q}(g_p) = \mathbb{Q}(\sqrt{p})$.
If $\left(\frac{-1}{p}\right) = -1$ (i.e., $p \equiv 3 \pmod 4$), then $g_p$ is purely imaginary and $g_p^2 = -p$, so $\mathbb{Q}(g_p) = \mathbb{Q}(\sqrt{-p})$.
In general, the unique quadratic extension is $\mathbb{Q}(\sqrt{p^*})$, where $p^* = (-1)^{(p-1)/2} p$.
Since $g_p$ is a linear combination of roots of unity with coefficients in $\mathbb{Q}$, $g_p \in \Phi_p$. Since $g_p^2 \in \mathbb{Q}$ and $g_p \notin \mathbb{Q}$, $\mathbb{Q}(g_p)$ is a quadratic extension of $\mathbb{Q}$ contained in $\Phi_p$. By uniqueness (part 2), this is the one.
\end{solution}

\begin{exercise}
1. Let $E/F$ be a finite Galois extension. Assume that the Galois group $\mathrm{Gal}(E/F)$ is generated by a single element $\sigma$. Let $x$ be an element of $E$ such that $\mathrm{tr}_{E/F}(x) = 0$. Show that there exists $y \in E$ such that $x = \sigma(y) - y$.

2. Let $F$ be a field of characteristic $p$, and let $E/F$ be a Galois extension of degree $p$. Show that there exists $x \in F$ such that $E \cong F[T]/(T^p - T - x)$.
\end{exercise}

\begin{solution}
1. This is the additive version of Hilbert's Theorem 90. Let $n = [E:F]$. The trace is $\mathrm{tr}_{E/F}(x) = x + \sigma(x) + \sigma^2(x) + \dots + \sigma^{n-1}(x)$.
By the linear independence of characters, there exists $z \in E$ such that $\mathrm{tr}_{E/F}(z) = \alpha \neq 0$. Scaling $z$, we can assume $\mathrm{tr}_{E/F}(z) = 1$.
Define $y$ as:
$$ y = \sum_{j=0}^{n-1} \left( \sum_{i=0}^j \sigma^i(x) \right) \sigma^j(z) $$
Let's check $\sigma(y) - y$:
$$ \sigma(y) = \sum_{j=0}^{n-1} \sigma\left( \sum_{i=0}^j \sigma^i(x) \right) \sigma^{j+1}(z) = \sum_{j=0}^{n-1} \left( \sum_{i=1}^{j+1} \sigma^i(x) \right) \sigma^{j+1}(z) $$
Since $\sigma^n = \mathrm{Id}$ and $\sigma^{j+1}$ is periodic, we can rearrange the sum.
Using the condition $\sum_{i=0}^{n-1} \sigma^i(x) = 0$, we find that $\sigma(y) - y = x$.

2. Since $[E:F] = p$ and $E/F$ is Galois, the Galois group $\mathrm{Gal}(E/F)$ is cyclic of order $p$. Let $\sigma$ be a generator.
In characteristic $p$, $\mathrm{tr}_{E/F}(1) = \sum_{i=0}^{p-1} \sigma^i(1) = \sum_{i=0}^{p-1} 1 = p \cdot 1 = 0$.
By part 1, there exists $\alpha \in E$ such that $\sigma(\alpha) - \alpha = 1$, i.e., $\sigma(\alpha) = \alpha + 1$.
Applying $\sigma$ repeatedly, we get $\sigma^i(\alpha) = \alpha + i$ for $i=0, 1, \dots, p-1$.
Since these $p$ values are distinct, the minimal polynomial of $\alpha$ over $F$ has degree $p$, so $E = F(\alpha)$.
Consider the element $x = \alpha^p - \alpha$. We have:
$\sigma(x) = \sigma(\alpha)^p - \sigma(\alpha) = (\alpha+1)^p - (\alpha+1) = \alpha^p + 1 - \alpha - 1 = \alpha^p - \alpha = x$.
Since $x$ is fixed by the generator $\sigma$ of the Galois group, $x \in F$.
The element $\alpha$ is a root of the polynomial $f(T) = T^p - T - x \in F[T]$.
Since $\alpha$ has degree $p$, $f(T)$ is the minimal polynomial of $\alpha$.
Thus $E = F(\alpha) \cong F[T]/(T^p - T - x)$.
\end{solution}

